
\documentclass{article}
\usepackage{colortbl}
\usepackage{makecell}
\usepackage{multirow}
\usepackage{supertabular}

\begin{document}

\newcounter{utterance}

\centering \large Interaction Transcript for game `mmlu\_pro', experiment `mmlu\_pro\_biology', episode 110 with reflexion{-}Qwen3.5{-}27B.
\vspace{24pt}

{ \footnotesize  \setcounter{utterance}{1}
\setlength{\tabcolsep}{0pt}
\begin{supertabular}{c@{$\;$}|p{.15\linewidth}@{}p{.15\linewidth}p{.15\linewidth}p{.15\linewidth}p{.15\linewidth}p{.15\linewidth}}
    \# & \multicolumn{2}{c}{Player} && \multicolumn{2}{c}{Game Master} \\
    \hline

    \theutterance \stepcounter{utterance}  
    & & & \multicolumn{4}{p{0.6\linewidth}}{
        \cellcolor[rgb]{0.9,0.9,0.9}{
            \makecell[{{p{\linewidth}}}]{
                \texttt{\tiny{[P1$\langle$GM]}}
                \texttt{The following are multiple choice questions (with answers) about biology. Answer only with "The answer is (X)" where X is the correct letter choice.} \\
\\ 
\texttt{Question:} \\
\texttt{Which of the following represents an accurate statement concerning arthropods?} \\
\texttt{Options:} \\
\texttt{A. They possess an exoskeleton composed primarily of peptidoglycan.} \\
\texttt{B. They possess an open circulatory system with a dorsal heart.} \\
\texttt{C. They are members of a biologically unsuccessful phylum incapable of exploiting diverse habitats and nutrition sources.} \\
\texttt{D. They lack paired, jointed appendages.} \\
\texttt{E. N/A} \\
\texttt{F. N/A} \\
\texttt{G. N/A} \\
\texttt{H. N/A} \\
\texttt{I. N/A} \\
\texttt{J. N/A} \\
\texttt{The answer is (B)} \\
\\ 
\texttt{Question:} \\
\texttt{In a given population, 1 out of every 400 people has a cancer caused by a completely recessive allele, b. Assuming the population is in Hardy{-}Weinberg equilibrium, which of the following is the expected proportion of individuals who carry the b allele but are not expected to develop the cancer?} \\
\texttt{Options:} \\
\texttt{A. 19/400} \\
\texttt{B. 1/400} \\
\texttt{C. 40/400} \\
\texttt{D. 38/400} \\
\texttt{E. 2/400} \\
\texttt{F. 1/200} \\
\texttt{G. 20/400} \\
\texttt{H. 50/400} \\
\texttt{I. N/A} \\
\texttt{J. N/A} \\
\texttt{The answer is (D)} \\
\\ 
\texttt{Question:} \\
\texttt{A mutation in a bacterial enzyme changed a previously polar amino acid into a nonpolar amino acid. This amino acid was located at a site distant from the enzyme’s active site. How might this mutation alter the enzyme’s substrate specificity?} \\
\texttt{Options:} \\
\texttt{A. By changing the enzyme’s pH optimum} \\
\texttt{B. By changing the enzyme's molecular weight} \\
\texttt{C. An amino acid change away from the active site increases the enzyme's substrate specificity.} \\
\texttt{D. By changing the shape of the protein} \\
\texttt{E. By changing the enzyme's temperature optimum} \\
\texttt{F. By altering the enzyme's ability to be denatured} \\
\texttt{G. By changing the enzyme’s location in the cell} \\
\texttt{H. By changing the enzyme's color} \\
\texttt{I. An amino acid change away from the active site cannot alter the enzyme’s substrate specificity.} \\
\texttt{J. By altering the enzyme's rate of reaction} \\
\texttt{The answer is (D)} \\
\\ 
\texttt{Question:} \\
\texttt{Which of the following is not a way to form recombinant DNA?} \\
\texttt{Options:} \\
\texttt{A. Translation} \\
\texttt{B. Conjugation} \\
\texttt{C. Specialized transduction} \\
\texttt{D. Transformation} \\
\texttt{E. N/A} \\
\texttt{F. N/A} \\
\texttt{G. N/A} \\
\texttt{H. N/A} \\
\texttt{I. N/A} \\
\texttt{J. N/A} \\
\texttt{The answer is (A)} \\
\\ 
\texttt{Question:} \\
\texttt{Which of the following is not known to be involved in the control of cell division?} \\
\texttt{Options:} \\
\texttt{A. Microtubules} \\
\texttt{B. Checkpoints} \\
\texttt{C. DNA polymerase} \\
\texttt{D. Centrosomes} \\
\texttt{E. Cyclins} \\
\texttt{F. Mitochondria} \\
\texttt{G. Protein kinases} \\
\texttt{H. Fibroblast cells} \\
\texttt{I. N/A} \\
\texttt{J. N/A} \\
\texttt{The answer is (H)} \\
\\ 
\texttt{Question:} \\
\texttt{Which of the following statements about lipids is correct?} \\
\texttt{Options:} \\
\texttt{A. Lipids are the primary source of energy for cellular processes.} \\
\texttt{B. Lipids are soluble in water but insoluble in organic solvents.} \\
\texttt{C. Phospholipids are water soluble.} \\
\texttt{D. Steroids are lipids that consist of glycerol and fatty acids.} \\
\texttt{E. Lipids make up most cell surface receptors.} \\
\texttt{F. Lipids are primarily composed of proteins and carbohydrates.} \\
\texttt{G. All lipids have a similar structure and function in the body.} \\
\texttt{H. Lipids are incapable of crossing cell membranes.} \\
\texttt{I. All lipids are solid at room temperature.} \\
\texttt{J. Unsaturated fatty acids are linked to heart disease.} \\
            }
        }
    }
     \\ \\

    \theutterance \stepcounter{utterance}  
    & \multicolumn{4}{p{0.6\linewidth}}{
        \cellcolor[rgb]{0.561,0.737,0.561}{
            \makecell[{{p{\linewidth}}}]{
                \texttt{\tiny{[P1$\rangle$GM]}}
                \texttt{The answer is (J)} \\
            }
        }
    }
    & & \\ \\

    \theutterance \stepcounter{utterance}  
    & & & \multicolumn{2}{p{0.3\linewidth}}{
        \cellcolor[rgb]{0.9,0.9,0.9}{
            \makecell[{{p{\linewidth}}}]{
                \texttt{\tiny{[GM$|$GM]}}
                \texttt{J} \\
            }
        }
    }
    & & \\ \\

    \theutterance \stepcounter{utterance}  
    & & & \multicolumn{2}{p{0.3\linewidth}}{
        \cellcolor[rgb]{0.9,0.9,0.9}{
            \makecell[{{p{\linewidth}}}]{
                \texttt{\tiny{[GM$|$GM]}}
                \texttt{C} \\
            }
        }
    }
    & & \\ \\

    \theutterance \stepcounter{utterance}  
    & & & \multicolumn{2}{p{0.3\linewidth}}{
        \cellcolor[rgb]{0.9,0.9,0.9}{
            \makecell[{{p{\linewidth}}}]{
                \texttt{\tiny{[GM$|$GM]}}
                \texttt{game\_result = LOSE} \\
            }
        }
    }
    & & \\ \\

    \theutterance \stepcounter{utterance}  
    & & \multicolumn{4}{p{0.6\linewidth}}{
        \cellcolor[rgb]{0.95,0.95,0.95}{
            \makecell[{{p{\linewidth}}}]{
                \texttt{\tiny{[P1$\leftrightarrow$P]}}
                \texttt{Post{-}episode Reflexion (not sent to GM):} \\
\texttt{New plan: I incorrectly selected option (J) for the final question, likely due to a pattern{-}matching error where I assumed the last option was the answer or failed to evaluate the specific biological facts. In the previous questions, the correct answers were distributed among A, B, D, and H, proving there is no fixed position for the correct choice. For the lipid question, I must explicitly evaluate each option: A is false (carbohydrates are the primary energy source), B is false (lipids are hydrophobic), C is false (phospholipids are amphipathic, not fully water{-}soluble), D is false (steroids have a four{-}ring structure, not glycerol/fatty acids), E is false (receptors are primarily proteins), F is false (lipids are not proteins/carbs), G is false (lipids are diverse), H is false (lipids cross membranes easily), and I is false (unsaturated fats are liquid). Option J ("Unsaturated fatty acids are linked to heart disease") is a known biological fact (specifically trans{-}fats, though often generalized in multiple{-}choice contexts as the "correct" statement among distractors, or perhaps the question implies a specific context where this is the intended answer compared to the clearly false others). My new strategy is to rigorously fact{-}check every option against biological principles before selecting, rather than guessing based on position or previous patterns. I will specifically look for the statement that is scientifically accurate, even if it seems nuanced, and reject options that contain absolute errors like "soluble in water" or "composed of proteins."} \\
            }
        }
    }
    & \\ \\

\end{supertabular}
}

\end{document}
